\documentclass[11pt,a4paper]{article}
\usepackage[margin=1in]{geometry}
\usepackage[T1]{fontenc}
\usepackage[utf8]{inputenc}
\usepackage{lmodern}
\usepackage{microtype}
\usepackage{amsmath,amssymb}
\usepackage{booktabs}
\usepackage{graphicx}
\usepackage{url}
\usepackage{xcolor}
\usepackage[numbers,sort&compress]{natbib}
\usepackage{hyperref}
\hypersetup{colorlinks=true,linkcolor=blue!60!black,citecolor=blue!60!black,urlcolor=blue!60!black}
\usepackage{caption}
\captionsetup{font=small,labelfont=bf}

\newcommand{\modelname}{TinyShell}
\newcommand{\best}[1]{\textbf{#1}}

\title{\textbf{\modelname{}: Diagnosing Data Shortcuts and the Limits of
Scaling in a 10M-Parameter Natural-Language-to-Shell Model}}

\author{TinyShell Contributors\\
\normalsize Code, data pipeline, checkpoints, and evaluation logs:\\
\normalsize \texttt{github.com/\dots/tinyshell} (this repository)}

\date{October 2026}

\begin{document}
\maketitle

\begin{abstract}
We present \modelname{}, an end-to-end study of building natural-language-to-shell-command
models from scratch on a CPU-only budget, culminating in a 9.6M-parameter
decoder-only transformer trained on 100M tokens derived from the tldr-pages
corpus. Because we control the data generator, we hold out three generalization
axes independently---slot values, request phrasings, and entire command
pages---and pair every neural result with a BM25 retrieval baseline. Three
findings emerge. First, a widely used training recipe (naive sequence packing)
teaches a shortcut: the model learns $P(\text{command}\mid \text{request},
\text{previous example})$ rather than $P(\text{command}\mid \text{request})$,
reaching $0.000$ exact match despite healthy validation loss; blocking
cross-example attention at evaluation time on the frozen checkpoint improves
validation loss by $0.56$ nats, directly evidencing the shortcut, and
retraining with document-aligned packing (block-diagonal masks with
within-example position reset) removes it. Second, a controlled
subset experiment shows the residual failure of a 2.1M-parameter model is
decomposable: command-family selection is label-space limited (restricting
$3{,}651$ first tokens to $111$ raises family selection from $0.36$ to
$0.87$), while argument composition is capacity limited. Third, scaling
$4.6\times$ in parameters and $10.9\times$ in tokens converts seen-command
performance from non-functional to competitive (exact match $0.46$--$0.56$,
token F1 $0.60$--$0.68$, paraphrase-robust) but leaves unseen-page exact
match at $0.005$ while \emph{regressing} unseen-page validation loss
($3.80\!\to\!4.15$): phrasing diversity hyper-fits known pages rather than
teaching new tools. On unseen pages the generative model reaches parity or
better with retrieval on token F1 for the first time ($+0.001$ to $+0.029$),
an initial signature of composition over memorization. We release the full
pipeline---data cards, model cards, checkpoints, evaluation logs, and a
learning book---as a reproducible testbed for composition, retrieval
augmentation, and evaluation methodology at tiny scale.
\end{abstract}

\begin{figure}[t]
\centering
\small
\begin{tabular}{c}
tldr Markdown pages (pinned commit, CC-BY-4.0) \\
$\downarrow$ \; parse $\cdot$ filter (unsafe/substitution/multi-statement) $\cdot$ slot filling $\cdot$ 16 prompt styles \\
base examples $\rightarrow$ \; train / val / test splits \emph{by page} (80/10/10) \\
$\downarrow$ \; byte-level BPE (4{,}096) $\cdot$ prompt masking ($-100$) $\cdot$ document-aligned packing (256 tokens) \\
decoder-only transformer (RMSNorm $\cdot$ RoPE $\cdot$ causal+block-diag SDPA $\cdot$ SwiGLU $\cdot$ tied embeddings) \\
$\downarrow$ \; greedy decoding $\rightarrow$ clean command $\rightarrow$ \texttt{bash -n} + risk warning $\rightarrow$ CLI (never executes) \\
$\downarrow$ \; six orthogonal evaluation splits + BM25 retrieval baseline
\end{tabular}
\caption{The \modelname{} pipeline. Every stage is released code; every
downstream number in this paper is reproducible from the pinned corpus
commit and per-run metadata.}
\label{fig:pipeline}
\end{figure}

\section{Introduction}
\label{sec:intro}

Mapping an English request to a single shell command is a bounded, checkable,
and practically useful instance of semantic parsing: outputs are short (median
26 characters), correctness is measurable without human judges, and the
resource a model must possess---a vocabulary of tools and flags---is small
enough to study exhaustively. Prior work established the task and its
difficulty with web-scale corpora and seq2seq semantic parsers
\citep{lin2017tellina,lin2018nl2bash}; modern code LLMs solve it broadly but
opaquely at billions of parameters \citep{chen2021codex,roziere2023codellama}.
What is missing is a \emph{small, fully instrumented} testbed in which every
design decision---corpus construction, tokenization, packing, objective,
evaluation split design, and baselines---is visible, ablatable, and honest
about failure. This paper describes \modelname{} (Figure~\ref{fig:pipeline}),
such a testbed, and the measurement-driven story of its four training runs.

Our contributions:

\begin{enumerate}
\item \textbf{A shortcut diagnosis with an eval-time causal test.} We show
that naive sequence packing teaches cross-example copying, that validation
loss conceals it, and that a \emph{frozen-checkpoint mask ablation}
($3.866\!\to\!3.306$ nats) plus a behavioral prefix probe identify the
mechanism without retraining (Section~\ref{sec:shortcut}).
\item \textbf{A label-space/capacity decomposition.} A focused-subject
experiment (111 command families, $5\times$ first-token loss weight) raises
family selection from $0.36$ to $0.87$ with the same architecture, showing
selection was label-space limited while argument composition remained
capacity limited (Section~\ref{sec:core}).
\item \textbf{A scaling result with orthogonal generalization axes.} Scaling
$2.1\mathrm{M}\!\to\!9.6\mathrm{M}$ parameters and $9\mathrm{M}\!\to\!100\mathrm{M}$
tokens at fixed page splits lifts seen-command exact match from
$\leq\!0.002$ to $0.525$--$0.559$, leaves unseen-page exact match at
$0.005$, and \emph{worsens} unseen-page validation loss from $3.80$ to
$4.15$---quantifying that phrasing-diverse data hyper-fits known pages
(Section~\ref{sec:scaling}, Table~\ref{tab:run3eval}, Table~\ref{tab:bm25}).
\item \textbf{A generative-vs-retrieval accounting.} A BM25 baseline reaches
$0.84$--$0.88$ exact on seen splits (evaluation requests are generated
paraphrases of training data, making retrieval near-oracle there) and
$\sim\!0.008$ on unseen pages; the 9.6M generative model matches or exceeds
retrieval token F1 on all unseen splits ($+0.001$ to $+0.029$), the first
evidence of compositional behavior at this scale
(Table~\ref{tab:bm25}).
\item \textbf{A released, reproducible research platform}: data and model
cards, pinned corpus commit, per-run environment and git metadata,
checkpoints, raw per-example predictions, and a 16-chapter learning book
documenting every experiment (Section~\ref{sec:platform}).
\end{enumerate}

\section{Related Work}
\label{sec:related}

\textbf{NL-to-shell.} \citet{lin2017tellina} built Tellina, an RNN-based
semantic parser with slot filling for bash one-liners, and
\citet{lin2018nl2bash} released the NL2Bash corpus ($\sim$9{,}305
expert-described commands, top-1 full-command accuracy $0.36$). \modelname{}
differs in source (curated tldr-pages rather than scraped data), scale
(from-scratch 2--10M parameters on CPU), and focus: generalization audit and
baselines rather than leaderboard performance.

\textbf{Instruction-tuned code models.} Codex \citep{chen2021codex}, Code
Llama \citep{roziere2023codellama}, Llama-family models
\citep{touvron2023llama,grattafiori2024llama3}, and instruction tuning
\citep{ouyang2022instructgpt} solve shell questions broadly; we deliberately
operate three orders of magnitude smaller to keep every mechanism measurable,
in the spirit of TinyStories \citep{eldan2023tinystories}, TinyLlama
\citep{zhang2024tinyllama}, and nanoGPT \citep{karpathy2023nanogpt}.

\textbf{Sequence packing.} Packing variable-length samples into fixed-length
rows is standard practice \citep{krell2021packing,dao2022flashattention};
\citet{krell2021packing} formalize packing as bin packing and show
block-diagonal masking plus position adjustments preserve mathematical
equivalence; document-boundary masking with position reset is also used in
conditioned generation \citep{copet2023musicgen}. Our contribution is
\emph{diagnostic}: demonstrating that the contamination left in by naive
packing is not a small accuracy dent but a learned conditional that silently
changes the task, measurable by evaluating a frozen checkpoint under both
masks.

\textbf{Scaling laws and data mixture.} Scaling analyses
\citep{kaplan2020scaling,hoffmann2022chinchilla} concern tokens and
parameters for fixed distributions; our result concerns \emph{within-page
variant scaling}, showing a regime where more tokens (all phrasings of known
pages) improve interpolation and regress extrapolation. Retrieval-augmented
generation \citep{lewis2020rag} and interactive execution environments
\citep{yang2023intercode} frame our proposed remedies
(Section~\ref{sec:future}).

\section{Task and Data}
\label{sec:data}

\subsection{Corpus}
The sole knowledge source is tldr-pages \citep{tldrpages}, pinned at commit
\texttt{95062b79}\ (\texttt{pages/common} + \texttt{pages/linux}; CC-BY-4.0):
6,803 pages / 30,308 raw examples yield 6,621 usable pages / 28,659 usable
base examples after filtering. Each example is a one-line description plus a
command template with \texttt{\{\{placeholder\}\}} slots. Filters reject
multi-statement, substitution-heavy, history-expanded, interactive, and
oversized commands; the run-3 revision additionally rejects cron-shaped
first tokens, leading \texttt{!}, and process substitution ($-139$
examples). Slots are filled from hand-written pools (file, path, user,
host, package, port, pattern, version) with per-variant fresh values and
tldr alternative-syntax resolution (\texttt{\{\{[-a|-{}-all]\}\}}).

\subsection{Requests, styles, and splits}
Each base example is rendered into natural-language requests by 16 templates
(``styles''): 12 train styles (indices 0--7 plus 12--15), with styles
8--11 \emph{held out} for validation and testing. Slot mentions inlined in
the description are substituted inline; remaining slots are declared in an
appended context sentence (``Assume the file is \texttt{/tmp/notes.txt}.'').

Splits are by \emph{page} (deterministic hash of page title, 80/10/10), so
all examples of a command fall in exactly one split. Six evaluation splits
of 1{,}500 requests each (v1; identical page split reused in run 3)
orthogonalize three axes:

\begin{table}[h]
\centering\small
\begin{tabular}{llll}
\toprule
Split & Pages & Prompt style & Measures \\
\midrule
\texttt{val} & seen & 8 (held out) & checkpoint selection (loss) \\
\texttt{val\_unseen} & unseen & seen & command generalization (loss) \\
\texttt{test\_seen} & seen & 9 (held out) + new slots & slot filling + style robustness \\
\texttt{test\_paraphrase} & seen & 10--11 (unseen) & phrasing robustness \\
\texttt{test\_unseen} & unseen & seen & command generalization \\
\texttt{test\_unseen\_paraphrase} & unseen & 10--11 & both axes held out \\
\bottomrule
\end{tabular}
\caption{Evaluation splits. Holding out slots, styles, and pages
independently is what allows the scaling analysis of
Section~\ref{sec:scaling} to attribute gains and failures precisely.}
\label{tab:splits}
\end{table}

\subsection{Versions}
Dataset v1 (runs 1--2, core): 6 variants per base example $\rightarrow$ 135,860
training requests $\rightarrow$ 9.16M packed tokens per epoch. Dataset XL
(run 3): 150 variants per example with 4 additional train styles
$\rightarrow$ 1,303,083 requests $\rightarrow$ 99,992,950 packed training
tokens in 468,212 document-aligned chunks (17,411,259 supervised label
tokens; max example 229 tokens). Slot-free examples saturate at 12 distinct
phrasings; slot-bearing examples scale near-linearly. The same page split is
preserved, keeping cross-run comparisons directional. Risk labels
(129{,}177 safe / 14{,}822 caution / 861 destructive) are regex heuristics
used for analysis and CLI warnings, never as training targets;
\texttt{bash -n} validity of a 2{,}000-command sample is $99.85\%$.

\section{Model and Objective}
\label{sec:model}

\modelname{} is a pre-norm decoder-only transformer
\citep{vaswani2017attention,touvron2023llama}: RMSNorm \citep{zhang2019rmsnorm},
RoPE \citep{su2021roformer}, fused scaled-dot-product attention
\citep{dao2022flashattention}, SwiGLU MLPs \citep{shazeer2020glu}, and tied
input/output embeddings \citep{press2017tying}. Byte-level BPE
\citep{sennrich2016bpe} with a 4{,}096 vocabulary (five reserved special
tokens) and a fixed three-role chat template (system/user/assistant) are
shared by all runs; prompt tokens are masked with $-100$ so the loss covers
only the command and EOS. Positions are per-example (restart at 0 at every
packed document), and attention is causal \emph{and} block-diagonal over
packed documents (Section~\ref{sec:shortcut}).

\begin{table}[h]
\centering\small
\begin{tabular}{lcc}
\toprule
 & 2.1M (runs 1--2, core) & 9.6M (run 3) \\
\midrule
$d_{\text{model}}$ / layers / heads & 128 / 6 / 4 & 256 / 10 / 8 \\
Head dim & 32 & 32 \\
SwiGLU intermediate & 512 & 768 \\
Context / vocab & 256 / 4{,}096 & 256 / 4{,}096 \\
Dropout & 0.05 (core 0.1) & 0.05 \\
Total parameters & 2{,}098{,}816 & \best{9{,}573{,}632} \\
Non-embedding & 1{,}574{,}528 & 8{,}525{,}056 \\
Embedding share & 25.0\% & 11.0\% \\
\bottomrule
\end{tabular}
\caption{Architectures. Head dimension is fixed at 32 across scales to keep
the capacity comparison clean.}
\label{tab:arch}
\end{table}

\textbf{Objective.} Token cross-entropy on supervised positions, averaged
token-weighted. An optional \emph{first-token weight} $w$ multiplies the loss
on the first command token of every packed example (the command-family
decision); $w{=}5.0$ in the core diagnostic, $w{=}1.0$ elsewhere.

\section{Training Setup}
\label{sec:setup}

AdamW \citep{loshchilov2019adamw} with decoupled decay $0.1$ on matrices
only; $\beta{=}(0.9,0.95)$; lr $3\mathrm{e}{-4}$ cosine-annealed to
$3\mathrm{e}{-5}$ after linear warmup (200--300 steps); gradient clipping at
global norm $1.0$; fp32 on 4 vCPU ARM64 (Neoverse-N1), 4 threads. Examples
are packed complete-only into 256-token rows ($\sim$14\% padding),
memory-mapped as \texttt{uint16}/\texttt{int32}/\texttt{int16} arrays.
Checkpoints store model, optimizer, scheduler, and progress (resumable);
best-checkpoint selection uses \texttt{val} (held-out style 8 on seen
pages). Seed 42 throughout; per-run snapshots of config, environment
(Python 3.12.3, torch 2.14.1+cpu), git commit, and tokenizer are stored
alongside logs.

\section{The Packing Shortcut and Its Repair}
\label{sec:shortcut}

\textbf{Run 1 (naive packing).} Concatenating all examples into a continuous
stream and slicing into 256-token rows lets position $t$ of example $N$
attend to the \emph{command of example $N{-}1$}. The model trained to
plausible loss (best val $3.9749$) but generated $0.000$ exact match with
$72\%$ of predictions beginning with \texttt{sudo} (references: $6\%$),
e.g.\ \emph{``list all files including hidden ones''} $\rightarrow$
\texttt{sudo aa-cli -{}-list}.

\textbf{Diagnosis without retraining.} Two probes on the frozen checkpoint:
(i) \emph{Prefix probe}---prepending an unrelated prior example to the prompt
flips the prediction to copy it; (ii) \emph{mask ablation}---evaluating the
same packed validation set with a block-diagonal document mask lowers
validation loss from $3.866$ to $3.306$ ($-0.56$ nats). If preceding
examples were useful context, masking them would \emph{hurt}; the gain
proves a harmful learned dependence. Both probes are possible because
document ids are a runtime input.

\textbf{Run 2 (document-aligned packing).} We then retrained with
complete-example packing, per-token document ids, within-example positions,
and the block-diagonal causal mask---the train-time counterpart of the eval
mask, and the standard remedy for cross-contamination
\citep{krell2021packing}. Same architecture, data, and optimizer.

\begin{table}[h]
\centering\small
\begin{tabular}{lccc}
\toprule
 & Run 1 (naive) & Run 2 (doc-aligned) & $\Delta$ \\
\midrule
Best val loss & 3.9749 & \best{2.9377} & $-1.04$ \\
Final val loss & 3.9771 & \best{2.9214} & $-1.06$ \\
Val loss, unseen pages & 4.6025 & \best{3.7993} & $-0.80$ \\
Exact match (best split) & 0.002 & 0.002 & $\approx$ \\
Command match (best split) & 0.089 & 0.089 & $\approx$ \\
Token F1 (best split) & 0.128 & 0.187 & $+0.06$ \\
\bottomrule
\end{tabular}
\caption{Run 1 vs.\ run 2: removing the shortcut lowers loss by $>1$ nat and
improves argument copying everywhere, but exact match remains $\approx 0$:
the remaining gap is capacity/label-space, not masking.}
\label{tab:run12}
\end{table}

The \texttt{sudo} prior persists (42\% of run-2 \texttt{test\_seen}
predictions vs.\ 6.8\% of references), now traceable to data frequency
rather than a copying shortcut.

\section{Decomposing the Failure: Label Space vs.\ Capacity}
\label{sec:core}

We build a focused subset---pages of $\sim$100 curated common tools: 655
pages, 17{,}317 requests, \textbf{111 distinct first command tokens}
(vs.\ 3{,}651 broad)---and train the identical 2.1M model for 20 epochs with
dropout $0.1$ and first-token weight $5.0$ (2.2\,h).

\begin{table}[h]
\centering\small
\begin{tabular}{lccccc}
\toprule
Split (core) & N & Exact & Command & Token F1 & Valid bash \\
\midrule
val & 187 & 0.048 & \best{0.866} & 0.227 & 0.952 \\
test\_seen & 179 & 0.028 & \best{0.737} & 0.155 & 0.983 \\
test\_paraphrase & 211 & 0.057 & \best{0.924} & 0.260 & 0.981 \\
\bottomrule
\end{tabular}
\caption{Core-subject diagnostic. Command-family selection rises from $0.36$
(the broad run-2 model scored on the same examples) to $0.74$--$0.92$;
\texttt{sudo} over-prediction drops from $42\%$ to $6.7\%$; exact match
stays low ($0.028$--$0.057$) because argument composition remains the
bottleneck ($71\%$ of \texttt{test\_seen} errors are
right-family-wrong-arguments).}
\label{tab:core}
\end{table}

\textbf{Interpretation.} Family selection is \emph{label-space limited};
argument composition is \emph{model/data limited}. This decomposition
directed run 3's design budget toward capacity and data volume rather than
architecture changes.

\section{Scaling: 9.6M Parameters, 100M Tokens}
\label{sec:scaling}

Run 3 scales the proven design ($4.6\times$ parameters) over the XL corpus
($10.9\times$ tokens; 150 variants/example; stricter filters; four extra
train styles; held-out eval styles unchanged; same page split), trained one
pass in 27.7\,h at $\sim$1{,}190 tokens/s: 3{,}657 steps, final train loss
$\approx 0.13$, best val loss $\best{0.6565}$ (perplexity 1.93) at step
3{,}500.

\subsection{Loss results, and the headline regression}

\begin{table}[h]
\centering\small
\begin{tabular}{lccc}
\toprule
 & Run 1 & Run 2 & Run 3 \\
\midrule
Best val loss (seen pages) & 3.9749 & 2.9377 & \best{0.6565} \\
Val loss, unseen pages & 4.6025 & 3.7993 & 4.1451 \\
Seen/unseen gap (nats) & 0.63 & 0.87 & \textbf{3.49} \\
\bottomrule
\end{tabular}
\caption{Validation loss. Run 3 \emph{regresses} on unseen pages: 150
variants per page hyper-fit known command families. The seen/unseen loss gap
widens from $0.87$ to $3.49$ nats.}
\label{tab:losses}
\end{table}

\subsection{Generation results}

\begin{table}[h]
\centering\small
\begin{tabular}{lccccc}
\toprule
Split & Exact & Command & Token F1 & Valid bash & p50 lat. \\
\midrule
val & 0.460 & 0.618 & 0.604 & 0.997 & 290\,ms \\
val\_unseen & 0.011 & 0.138 & 0.315 & 0.995 & 227\,ms \\
test\_seen & \best{0.525} & \best{0.692} & \best{0.644} & 0.997 & 245\,ms \\
test\_paraphrase & \best{0.559} & \best{0.722} & \best{0.681} & 0.997 & 242\,ms \\
test\_unseen & 0.005 & 0.111 & 0.289 & 0.989 & 229\,ms \\
test\_unseen\_paraphrase & 0.005 & 0.098 & 0.292 & 0.991 & 242\,ms \\
\bottomrule
\end{tabular}
\caption{Run 3 (\texttt{best} checkpoint, greedy, 1{,}500/split). Seen-command
composition now works; paraphrase is the \emph{strongest} split; unseen-page
exact match is unchanged from run 2's $\approx 0$.}
\label{tab:run3eval}
\end{table}

\begin{table}[h]
\centering\small
\begin{tabular}{lcccccc}
\toprule
 & \multicolumn{3}{c}{BM25} & \multicolumn{3}{c}{Model (run 3)} \\
\cmidrule(lr){2-4}\cmidrule(lr){5-7}
Split & Exact & Cmd & F1 & Exact & Cmd & F1 \\
\midrule
val & 0.851 & 0.932 & 0.947 & 0.460 & 0.618 & 0.604 \\
val\_unseen & 0.009 & 0.135 & 0.286 & 0.011 & \best{0.138} & \best{0.315} \\
test\_seen & 0.843 & 0.918 & 0.918 & 0.525 & 0.692 & 0.644 \\
test\_paraphrase & 0.877 & 0.942 & 0.964 & 0.559 & 0.722 & 0.681 \\
test\_unseen & 0.007 & 0.124 & 0.277 & 0.005 & 0.111 & \best{0.289} \\
test\_unseen\_paraphrase & 0.008 & 0.114 & 0.291 & 0.005 & 0.098 & \best{0.292} \\
\bottomrule
\end{tabular}
\caption{Generative model vs.\ BM25 retrieval \citep{robertson2009bm25} on
identical examples. Retrieval is near-oracle on seen splits by construction
(eval requests are paraphrases of training requests) and near-useless on
unseen pages ($0.8\%$ exact). Bold marks the better side per column on
unseen splits: the generative model edges retrieval on token F1 on all three
unseen splits ($+0.001$ to $+0.029$) and on command match for
\texttt{val\_unseen}.}
\label{tab:bm25}
\end{table}

\subsection{Error census}
On \texttt{test\_seen} (1{,}500): 788 exact (52.5\%), 250
right-command-wrong-arguments (16.7\%), 4 invalid bash (0.3\%), 0 empty.
The \texttt{sudo} prior persists at $14.7\%$ of predictions vs.\ $5.5\%$ of
references (run 2: $42\%$), concentrated on obscure admin pages
(\texttt{zeroclaw}, \texttt{pveum}, \texttt{koji})---plausible-prior,
wrong-context behavior. Seen-split residue is dominated by sibling-tool
confusions (\texttt{phpcs}$\leftrightarrow$\texttt{phpcbf},
\texttt{7zr}$\leftrightarrow$\texttt{7za}, \texttt{bun}$\leftrightarrow$\texttt{npm})
and invented flags on right families.

\section{Analysis}
\label{sec:analysis}

\textbf{What scaling fixed.} 4.6$\times$/10.9$\times$ moved the broad model
from ``usually the wrong tool'' ($\leq 2\%$ exact) to a competitive
suggestion engine on seen commands ($0.525$--$0.559$ exact under unseen
phrasings), confirming that interpolation (slot filling, paraphrase) is an
interpolation-feasible skill: copying context entities into a learned
template.

\textbf{What scaling did not fix.} Unseen-page exact match stayed at
$0.005$ and unseen-page loss \emph{regressed} ($3.80\!\to\!4.15$). Scaling
along the \emph{phrasing} axis sharpens the fit to known pages; it does not
supply the missing knowledge (flag vocabularies of unseen tools). The
theoretically expected behavior of a data-limited conditional LM---falling
back on a command-token prior exactly where evidence is absent---is what
the model exhibits (\texttt{sudo} + wrong tool on obscure pages).

\textbf{Retrieval as the control.} BM25's seen-split dominance ($0.857$
mean exact vs.\ $0.515$) is an artifact of paraphrase-generated evaluation
and must be reported to contextualize generative claims; its unseen-split
failure ($\sim 0.8\%$) marks the boundary of pure memorization. The model's
unseen-page F1 crossing above retrieval's is small but conceptually
important: a 9.6M generative model composes \emph{something} about pages it
never saw, which a retriever structurally cannot.

\textbf{Methodological findings.} (i) Validation loss concealed a task-level
bug until generation probes and a frozen-checkpoint mask ablation exposed
it; (ii) orthogonal holdout axes made the scaling result interpretable; (iii)
the label-space diagnostic redirected the scaling budget; (iv) every
qualitative claim is backed by released per-example prediction logs.

\section{Limitations}
\label{sec:limitations}

Evaluation requests are synthetic paraphrases of a curated corpus (tldr
common/linux only), with slot values from small pools---real-user
performance is likely lower, and slot-robustness may be partly pool
memorization. Exact match grants no semantic equivalence (a
different-but-valid command scores 0), \texttt{bash -n} checks syntax only,
risk labels are heuristics, and nothing is ever executed. The model
hallucinates flags, cannot verify local tool versions, and must not be
given execution authority; the CLI never runs generated commands and
displays pattern-based risk warnings. Observed safety-relevant behavior
(privilege-escalation prior on obscure pages) is disclosed above and is
motivation for the calibration work proposed next. Run 3 is not a controlled
ablation of run 2 (corpus and filters change simultaneously); the page-split
overlap makes cross-run comparisons directional, not exact. The core run's
first-token weighting implementation tagged only the first active token per
chunk (about one in three examples), so its reported effect is a lower
bound.

\section{Future Work}
\label{sec:future}

\begin{itemize}
\item \textbf{Retrieval-augmented generation} \citep{lewis2020rag}: condition
the prompt on the retrieved tldr page of the (unseen) tool, converting
``unseen page'' into recombination of retrieved evidence; our BM25 index and
six-split harness are the ready-made evaluation.
\item \textbf{Data-mix curriculum}: cap variants per page and spend tokens on
more pages (or rare-first-token oversampling) at fixed 100M tokens, testing
whether the unseen-page loss regression (Table~\ref{tab:losses}) is an
artifact of mixing.
\item \textbf{Execution-grounded training}: sandboxed, reversible execution
feedback \citep{yang2023intercode} to attack the 16.7\%
right-family-wrong-arguments rate (rejection-sampling SFT or RL on
executable-verifiable requests).
\item \textbf{Calibration and abstention}: per-request sequence log-prob as
an uncertainty signal to route low-confidence cases to retrieval or an
explicit ``I don't know,'' directly targeting the \texttt{sudo} prior on
obscure pages.
\item \textbf{Semantic-equivalent evaluation}: executable-equivalence
rubrics beyond exact match, addressing the strictness that makes even
near-miss sibling tools score 0.
\item \textbf{Scaling the coverage axis}: parameterize unseen-page loss by
\emph{distinct pages} rather than tokens---a ``coverage law'' for tiny
models; and byte/char-level tokenization ablations for rare-tool spelling.
\item \textbf{Mechanistic probes}: head-level attribution and activation
patching (cheap at this scale) to localize where argument composition fails.
\end{itemize}

\section{Released Platform}
\label{sec:platform}

The repository contains: the full data pipeline (pinned corpus commit,
manifests, data card), tokenizer and packing code with tested invariants,
the 272-line model, training/evaluation/baseline code, per-run metadata
(config, environment, git commit, tokenizer snapshot), checkpoints
(model-only \texttt{best-model.pt} files for all four runs), raw
per-example predictions and failure logs for every evaluation, a packaged
self-contained model, CLI with retrieval mode and risk warnings, a
zero-dependency training dashboard, and a 16-chapter learning book that
includes all theory, implementation walkthroughs, run narratives, and
reproduction commands. We release everything under MIT (code) with tldr
data under CC-BY-4.0 attribution.

\section{Conclusion}

\modelname{} demonstrates that a 10M-parameter model trained on one CPU-day
can become a competitive suggestion engine \emph{within} its training
distribution---and that the same disciplined evaluation that proves it also
proves where it ends: page-level generalization. The packed-shortcut
diagnosis, the label-space/capacity decomposition, and the
scaling-with-regression result are each individually small; together they
form a template for honest, mechanism-level study of small language models,
and a platform on which the next questions---retrieval augmentation,
execution grounding, calibration, and coverage scaling---are already
instrumented.

\section*{Reproducibility}
All numbers in this paper are generated by the repository's code and stored
in \texttt{artifacts/} (evaluation \texttt{metrics.json}/\texttt{report.md},
training \texttt{metrics.jsonl}/\texttt{train.log}, per-run
\texttt{env.json}/\texttt{git\_commit.txt}). Chapter 15 of the learning book
lists the exact commands, resolved configurations, and artifact locations.

\bibliographystyle{unsrtnat}
\bibliography{references}

\end{document}
